\section{Unipotent representations of p-adic groups} \label{sec:padic_groups}

Let $F$ be a finite extension of $\QQ_p$ with residue field $\FF_q$, and let $\mathcal{G}$ be a connected simple algebraic $F$-group split over $F$ and write $G=\mathcal{G}(F)$. Let $\mathcal{T}$ be a $F$-split maximal torus of $\mathcal{G}$ and let $T=\mathcal{T}(F)$ be a maximal torus of $G$. We assume throughout that $\mathcal{G}$ is in addition of simply-connected type. 
%Let $G$ be a connected simple group defined over a finite extension $F$ of $\QQ_p$ with residue field $\FF_q$. We further assume that $G$ is split with split maximal torus $T$ and of simply connected type. 

Let $\Irr(G)$ be the set of irreducible smooth representations of $G$ and $\Irr_{\un}(G)\subseteq\Irr(G)$ the subset of unipotent representations. Under the Local Langlands correspondence, there is a bijection
\begin{align*}
    \mathrm{LLC}_{\un}:\left\{
    \begin{array}{c}
        \text{Unipotent representations} \\
        \text{of } G(F)
    \end{array}
    \right\}_{/\cong}
    &\longleftrightarrow
    \left\{
    \begin{array}{c}
        \text{Pairs }(x,\phi):\\
        x\in G^\vee,\ \phi\in\widehat{A_x}       
    \end{array}
    \right\}_{/G^\vee},
\end{align*}

By the work of Lusztig, unipotent representations of $G$ can also be characterized directly in terms of the topology of $G$ and the representation theory of finite groups of Lie type. For each facet $c$ in the Bruhat--Tits building $\mathcal{B}(G,F)$, consider the parahoric subgroup $K_c:=\Stab_G(c)$. Since, by assumption, $G$ is simply connected, $K_c$ is a connected open compact subgroup of $G$ fitting in a short exact sequence
$$1\longrightarrow U_c\longrightarrow K_c\longrightarrow \overline{K}_c\longrightarrow 1,$$
where $U_c$ is the pro-unipotent radical of $K_c$ and $\overline{K}_c$ is the group of $\FF_q$-rational points of a connected reductive group $\mathbf{\overline{K}_c}$. Since $K_{g\cdot c}=gK_cg^{-1}$ for all facets $c$ of $\mathcal{B}(G,F)$ and $g\in G$, conjugacy classes of parahoric subgroups are in bijection with $G$-orbits of facets in $\mathcal{B}(G)$. Furthermore, after fixing a set of simple affine roots $S_{\aff}$ of $G$ with corresponding fundamental alcove $\mathcal{C}_0$, there is a bijection
\begin{align*}
    \{\text{Subsets }J\subsetneq S_{\aff}\}&\longleftrightarrow\{G\text{-orbits of facets in }\mathcal{B}(G)\}\longleftrightarrow\{\text{Conj. classes of parahoric subgroups}\}\\
    J&\longmapsto c_J=\{x\in\overline{\mathcal{C}_0}\ |\ \langle\alpha,x\rangle\in\ZZ\iff\alpha\in J,\text{ for any }\alpha\in S_{\aff}\}\longmapsto K_{c_J}
\end{align*}
satisfying that $J_1\subseteq J_2$ if and only if $c_{J_2}\subset \overline{c_{J_1}}$ if and only if $K_{c_{J_1}}\subseteq K_{c_{J_2}}$. Throughout, we denote by $(K_J,U_J,\overline{K}_J)$ the parahoric subgroup corresponding to the facet $c_J$. In particular, maximal parahoric subgroups correspond to facets of minimal dimension and maximal subsets $J\subsetneq S_{\aff}$. When $G$ is simply connected, we recall that the set of maximal parahoric subgroups coincides with the set of maximal open compact groups. We define $$\max(G):=\{K_x\ |\ \{x\}\subset\mathcal{B}(G)\text{ facet of minimal dimenson}\}$$
to be the set containing each orbit of maximal parahoric subgroups.

Next, let $(\pi,V)$ a smooth admissible complex representation of $G$ and let $(K,U,\overline{K})$ be a parahoric subgroup. The space $V^{U_K}$ of fixed points under the pro-unipotent radical is naturally a representation of $\overline{K}=K/U_K$. This induces the \textit{parahoric restriction functor}
\begin{equation}\label{eqn:par_restriction}
    \res_K:R(G)\longrightarrow R(\overline{K}),\quad V\longmapsto V^{U_K},\quad\text{for all }V\in\Irr(G),
\end{equation}
where $R(G)$ and $R(\overline{K})$ are the complexification of the Grothendieck group of $\Irr(G)$ and $\Irr(\overline{K})$, respectively. If $(\sigma,E)$ is a cuspidal representation of $\overline{K}$, we define
$$\Irr(G,K,E)=\{(\pi,V)\in\Irr(G)\ |\ \text{ the }\overline{K}\text{-module $V^{U_K}$ contains }E\}.$$
\begin{definition}
    We say that an irreducible representation $(\pi,V)$ of $G$ is \textit{unipotent} if there is a parahoric subgroup $(K,U_K,\overline{K})$ such that $V^{U_K}$ contains a cuspidal unipotent representation of $\overline{K}$; that is, if $(\pi,V)\in\Irr(G,K,E)$ for some pair $(K,E)$ where $E$ is unipotent.
\end{definition}
Therefore, the set of unipotent representations, up to isomorphism, is
$$\Irr_{\un}(G)=\bigcup_{\substack{J\subsetneq S_{\aff}\\ E \text{ cusp. unip. $\overline{K}_J$-rep}}}\Irr(G,K_J,E).$$
Each set $\Irr(G,K,E)$ generates a distinct \textit{Bernstein block} of representations of $G$. If $\mathcal{I}$ is an Iwahori subgroup, $\mathcal{I}/U_\mathcal{I}\cong(\FF_q^\times)^r$ and the only unipotent representation is the trivial character. Thus, Iwahori-spherical representations are all unipotent. 

Moreover, unipotent representations behave well under the parahoric restriction functors $\res^K$. The first important result is due to Moy and Prasad. A similar result holds for non-unipotent representations, but it is more complicated. 
\begin{proposition}\label{thm:unip_restriction}\cite{Reeder2000}*{Lemma 5.1}
    Assume $G$ is simply connected. Suppose $I\subsetneq S_{\aff}$ and that $V^{U_I}$ contains a cuspidal unipotent representation $\sigma$ of $\overline{K}_I$. If $J\subsetneq S_{\aff}$ with $V^{U_J}\neq 0$, and $J$ is minimal with respect to this property $I=J$.
\end{proposition}
The following is a direct consequence of the above result.
\begin{cor}\label{cor:unipotentblock_disjoint}
    For any two pairs $(K,E)$ and $(K',E')$ of a parahoric subgroup and a cuspidal unipotent representation of the reductive quotient, $\Irr(G,K,E)$ and $\Irr(G,K',E')$ are either disjoint or equal.
\end{cor}
Proposition \ref{thm:unip_restriction} together with a standard Harish--Chandra style argument yields the following corollary.
\begin{cor}
    Let $(\pi,V)$ be a unipotent representation of $G$ and let $(K,U_K,\overline{K})$ be a parahoric subgroup. Then the $\overline{K}$-irreducible components of $V^{U_K}$ are all unipotent.
\end{cor}

Consequently, for any parahoric subgroup $K$, the parahoric restriction functor $\res^K$ naturally restricts to the \textit{unipotent parahoric restriction function}
\begin{equation*}
    \res^K_{\un}:R_{\un}(G)\longrightarrow R_{\un}(\overline{K}),\quad V\longmapsto V^{U_K},
\end{equation*}
where $R_{\un}(G)$ and $R_{\un}(\overline{K})$ are the complexification of the Grothendieck groups spanned by $\Irr_{\un}(G)$ and $\Irr_{\un}(\overline{K})$, respectively. Finally, we let
$$\res^{\cpt}_{\un}:=\bigoplus_{K\in\max(G)}\res^K_{\un}:R_{\un}(G)\longrightarrow \bigoplus_{K\in\max(G)}R_{\un}(\overline{K}).$$

\subsection{The Unipotent Langlands classification}\label{subsec:ULLC}

The Local Langlands provides a beautiful classification of the unipotent representations of $G$ in terms of the geometry of the dual group $G^\vee$ of $G$, of adjoint type. This result is widely known as Lusztig's arithmetic-geometric correspondence.
\begin{theorem}[The arithmetic-geometric correspondence \cite{Lusztig1995}] There is a natural bijection
    \begin{align*}
        G^\vee\backslash\{(x,\phi)\ |\ x\in G^\vee,\phi\in\widehat{A_x}\}&\longleftrightarrow \Irr_{\un}(G)\\
        (x,\phi)&\longmapsto \pi(x,\phi).
    \end{align*}
    In particular, there is an isomorphism
    \begin{equation}\label{eqn:arithmetic-geometric}
        \LLC_{\un}:\bigoplus_{x\in G^\vee}R(A_x)\longrightarrow R_{\un}(G),
    \end{equation}
    where the sum on the left hand side is taken over conjugacy classes of $G^\vee$ and $R(A_x)$ is the space of virtual characters of $A_x$.
\end{theorem}

Given a Langlands parameter $(x,\phi)$, it is a natural question to ask inside which unipotent block does the representation $\pi(x,\phi)$ lies in. This can be answered using the generalized Springer correspondence, which we will not revise in detail here (see \cite{CiubotaruMasonBrownOkada2025}*{\S4.3}). We only state a fundamental result giving a geometric characterization for a representation $\pi(x,\phi)$ to be Iwahori-spherical, and thus lying in the principal block generated by $\Irr(G,\mathcal{I},\mathbf{1})$. Let $\mathcal{B}^x$ the partial flag variety containing all Borel subgroups of $G^\vee$ containing $x$. This variety carries a natural action of $A_x$ given by conjugation. In particular, the cohomology complex $H^*(\mathcal{B}^x)$ inherits an action of $A_x$ too.

\begin{theorem}\cite{KazhdanLusztig1987}*{Theorem 7.12}
    The representation $\pi(x,\phi)$ is Iwahori-spherical if and only if $H^*(\mathcal{B}^x)^\phi\neq 0$. In particular, for any $x\in G^\vee$, the representation $\pi(x,\mathbf{1})$ is Iwahori-spherical.
\end{theorem}
Consequently, unipotent representations can also be described as the smallest class of representations of $G$ containing Iwahori-spherical representations and closed under L-packets.

In later sections, we display parahoric restriction tables for representations of $G$ parametrized by several elements $x\in G^\vee$. To do this, we use the following standard notation for the Langlands parameters as in \cite{Reeder2000}, which we now explain. 
\begin{notation}\label{notation:Langlandsparam}
    Given a representation $\pi(x,\phi)$, let $x=us$ be the Jordan decomposition. Then $u$ is a unipotent element of the connected pseudo-Levi subgroup $Z_{G^\vee}(s)^0$, and $\phi$ is an irreducible representation of the component group $A_{us}$. The representation $\pi(x,\phi)$ will be labelled by the type of unipotent element $u$ is inside $Z_{G^\vee}(s)^0$, together with the representation $\phi$. For example, if $G^\vee=F_4(\CC)$ and $u\in G^\vee$ is a unipotent element of type $F_4(a_3)$, there is a semisimple element $s\in Z_{G^\vee}(u)$ of order $2$ such that $Z_{G^\vee}(s)$ is connected of type $C_3\times A_1$. The unipotent element $u$ is then subregular in the first component and regular in the second, labelled by $C_3(42)\times A_1$, and $\widehat{A_{us}}=C_2\times C_2$. Thus, the four representations of $F_4(F)$ in the $L$-packet parametrized by $x=su\in F_4(\CC)$ are labelled $[C_3(42)A_1,\pm\pm]$. %Thus, each representation has a name of the form $[u\in Z_{G^\vee}(s)^0,\phi]$ where $\phi$ is an irreducible representation of $A_{us}$.
\end{notation}


\subsection{Chevalley involutions and the dual nonabelian Fourier transform}\label{subsec:Fourier_transform}

In this section we define a modified version of the \textit{dual elliptic nonabelian Fourier transform}, constructed in \cite{AubertCiubotaruRomano2024}*{\S9} \textcolor{red}{Add citations}. The missing ingredient to the conjecture is the role played by a Chevalley involution on the Langlands parameters, which we now describe. As usual, we let $G^\vee$ be the complex dual group of $G$.

\begin{definition}[\textcolor{red}{cite this better}]
    An automorphism $C:G^\vee\to G^\vee$ is called a \textit{Chevalley involution} if there is some maximal torus $T^\vee$ of $G^\vee$ such that $C(t)=t^{-1}$ for all $t\in T^\vee$.
\end{definition}
Let $\Aut(G^\vee)$ be the set of automorphisms of $G^\vee$, $\Int(G^\vee)$ the subset of inner automorphisms and $\mathrm{Out}(G^\vee)=\Aut(G^\vee)/\Int(G^\vee)$. We denote inner automorphisms by $\inn(x):g\mapsto xgx^{-1}$. Using based root data and pinnings, one can show that any Chevalley involution $C:G^\vee\to G^\vee$ exists, satisfies $C^2=1$ and for every semisimple element $g\in G$, $C(g)$ is conjugate to $g$. In addition, if $C_1$ and $C_2$ are Chevalley involutions with respect to $T_1^\vee$ and $T_2^\vee$, then they are conjugate by an inner automorphism of of $G^\vee$ (\textcolor{red}{cite here!!}). Thus, Chevalley involutions all lie in a distinguished class of $\mathrm{Out}(G^\vee)$ of order dividing $2$. The following result characterizes the order of this class depending on the type of $G^\vee$. We insist that $G^\vee$ is assumed to be simple.

\begin{lemma}
    The Chevalley involutions $C:G^\vee\to G^\vee$ are inner if and only if $-1$ is in the Weyl group, in which case, $C=\inn(n_0)$, for some $n_0\in N_{G^\vee}(T^\vee)$ representing the longest element $w_0=-1$. In particular, Chevalley involutions are inner for types $A_1,B_n,C_n,D_{2k},E_7,E_8,F_4$ and $G_2$, and outer for the other types.
\end{lemma}


The \textit{dual nonabelian Fourier transform} is a conjectural involution $\FT^\vee:R_{\un}(G)\to R_{\un}(G)$ that lifts Lusztig's non-abelian Fourier transform and is compatible with the parahoric restriction maps $\res^K_{\un}$. While such a map has not been constructed in the literature on the infinite dimensional space $R_{\un}(G)$, we consider the finite-dimensional quotient $\overline{R}_{\un}(G)$ of $R_{\un}(G)$ by the radical of the Euler--Poincare pairing

\begin{equation*}
    EP_G(\ ,\ ):R_{\un}(G)\times R_{\un}(G)\longrightarrow\CC,\quad (V,V')\longmapsto\sum_{i\geq0}(-1)^i\Ext^i(V,V'),\text{ where }V,V'\in\Irr_{\un}(G),
\end{equation*}
and extended by linearity, and where $\Ext^i(V,V')$ are calculated in the category of smooth complex $G$-representations. Thus, the space $\overline{R}_{\un}(G)$ carries a natural non-degenerate pairing $EP_G(\ ,\ )$. Similarly, for each $x\in G^\vee$ the space of virtual representations $R(A_x)$ carries a natural elliptic pairing $(\ ,\ )^{A_x}_{\mathrm{ell}}$ (see \cite{AubertCiubotaruRomano2024}*{\S8.1} for the definition), and we let $\overline{R}(A_x)$ be the quotient of $R(A_x)$ by the radical of the elliptic pairing. The arithmetic-geometric \eqref{eqn:arithmetic-geometric} induces an isometric isomorphism on the elliptic representation spaces
\begin{equation}\label{eqn:LLC_unipotent_elliptic}
    \overline{\LLC}_{\un}:\bigoplus_{x\in G^\vee}\overline{R}(A_x)\longrightarrow\overline{R}_{\un}(G),
\end{equation}
where the sum on the left is again over conjugacy classes of elements in $G^\vee$. The left hand side of \eqref{eqn:LLC_unipotent_elliptic} can be further decomposed into
$$\bigoplus_{x\in G^\vee}\overline{R}(A_x)=\bigoplus_{\substack{u\in G^\vee\\\mathrm{unipotent}}}\ \bigoplus_{\substack{s\in Z_{G^\vee}(u)\\\mathrm{semisimple}}}\overline{R}(A_{us}).$$
Next, fix some unipotent element $u\in G^\vee$ and let $\Gamma_u$ be the reductive part of the centralizer $Z_{G^\vee}(u)$. We define the set of elliptic pairs associated to $u$ to be
$$\mathcal{Y}(\Gamma_u)_{\mathrm{ell}}=\{(s,h)\in\Gamma_u\times\Gamma_u\ |\ s,h\text{ semisimple, }sh=hs, Z_{\Gamma_u}(s,h)\text{ is finite}\}.$$
It is not hard to show that this set is always finite (see \cite{AubertCiubotaruRomano2024}*{Lemma 8.1}). For every $(s,h)\in\mathcal{Y}(\Gamma_u)$, define 
\begin{equation}\label{eqn:virtual_ellipticpair}
    \Pi(s,h)=\sum_{\phi\in\widehat{A_{\Gamma_u}(s)}}\phi(\overline{h})\phi\in R(A_{\Gamma_u}(s)),
\end{equation}
where $\overline{h}$ is the image of $h$ in $A_{\Gamma_u}(s)$. Moreover, let $\overline{\Pi}(s,h)$ denote the image in $\overline{R}(A_{\Gamma_u}(s))$ and let $\CC[\mathcal{Y}(\Gamma_u)_{\mathrm{ell}}]^{\Gamma_u}\cong\CC[\Gamma_u\backslash\mathcal{Y}(\Gamma_u)_{\mathrm{ell}}]$ be the space of $\Gamma_u$-invariant functions $\mathcal{Y}(\Gamma_u)$. Let $\mathbf{1}_{[(s,h)]}$ be the indicator function of the $\Gamma_u$-orbit of $(s,h)$. 
\begin{proposition}\label{prop:elliptic_pairs}\cite{AubertCiubotaruRomano2024}*{Proposition 8.3} Let $u\in G^\vee$ be unipotent. Then the correspondence $\mathbf{1}_{[(s,h)]}\mapsto\overline{\Pi}(s,h)$ induces an isomorphism
    $$\CC[\mathcal{Y}(\Gamma_u)_{\mathrm{ell}}]^{\Gamma_u}\longrightarrow\bigoplus_{\substack{s\in Z_{G^\vee}(u)\\\mathrm{semisimple}}}\overline{R}(A_{us})$$
\end{proposition}
Finally, for each $[(s,h)]\in\Gamma_u\backslash\mathcal{Y}(\Gamma_u)_{\mathrm{ell}}$, define the virtual combinations
\begin{equation}
    \Pi(u,s,h):=\sum_{\phi\in\widehat{A_{us}}}\phi(h)\pi(us,\phi)\in R_{\un}(G),
\end{equation}
and let $\overline{\Pi}(u,s,h)$ be the corresponding images in $\overline{R}_{\un}(G)$.
\begin{definition}
    The elliptic unipotent representation space $R_{\un,\mathrm{ell}}(G)$ is the finite-dimensional subspace of $R_{\un}(G)$ spanned by virtual characters 
    $\{\Pi(u,s,h)\ |\ u\text{ unipotent,  }(s,h)\in\mathcal{Y}(\Gamma_u)_{\mathrm{ell}}\}.$
\end{definition}
Thus, the isomorphism \eqref{eqn:LLC_unipotent_elliptic} together with Proposition \ref{prop:elliptic_pairs} yields an isometric isomorphism
\begin{equation}
    R_{\un,\mathrm{ell}}(G)\longrightarrow\overline{R}_{\un}(G),\quad\Pi(u,s,h)\mapsto\overline{\Pi}(u,s,h)
\end{equation}
We now define the dual elliptic nonabelian Fourier transform on $R_{\un,\mathrm{ell}}(G)$. 

\begin{definition}
    The \textit{dual elliptic nonabelian Fourier transform} is the involutive linear map $\FT^\vee_{\mathrm{ell}}:R_{\un,\mathrm{ell}}\rightarrow R_{\un,\mathrm{ell}}$ defined by
    \begin{equation*}
        \FT^\vee_{\mathrm{ell}}(\Pi(u,s,h))=\Pi(C(u),C(h),C(s)),\quad u\in G^\vee\text{ unipotent, }(s,h)\in\Gamma_u\backslash\mathcal{Y}(\Gamma_u)_{\mathrm{ell}},
    \end{equation*}
    where $C:G^\vee\to G^\vee$ is a Chevalley involution of the complex dual group.
\end{definition}
We note that $\FT^\vee_{\ell}$ is a well-defined involution. A Chevalley involution induced an isomorphism $C:\Gamma_u\to\Gamma_{C(u)}$ and therefore $(C(h),C(s))\in\Gamma_{C(u)}\backslash\mathcal{Y}(\Gamma_{C(u)})_{\mathrm{ell}}$. Moreover, any two Chevalley involutions differ by an inner automorphism, so $\FT^\vee_{\mathrm{ell}}$ is independent on the choice of Chevalley involution. Finally, 
$$(\FT^\vee_{\mathrm{ell}})^2(\Pi(u,s,h))=\Pi(C^2(u),C^2(s),C^2(h))=\Pi(u,s,h).$$
%We are now ready to state the main conjecture for elliptic unipotent representations.
\begin{conjecture}\cite{AubertCiubotaruRomano2024}*{Conjecture 1.2}\label{conj:main2}
    Let $G$ be a split semisimple $p$-adic group of simply connected type. Then the following diagram commutes up to roots of unity,
    \begin{equation}\label{eqn:diagram}
        \begin{tikzcd}
        R_{\un,\text{ell}}(G) \arrow[r, "\FT^\vee_{\text{ell}}"] \arrow[d, "\res_{\un}^{\cpt}"] & R_{\un,\text{ell}}(G) \arrow[d, "\res_{\un}^{\cpt}"] \\
        \bigoplus_{K\in\max(G)}R_{\un}(\overline{K}) \arrow[r, "\FT^{\cpt}_{\un}"] & \bigoplus_{K\in\max(G)}R_{\un}(\overline{K}),
        \end{tikzcd}  
    \end{equation}
    where $\FT^{\cpt}_{\un}:=\bigoplus_{K\in\max(G)}\FT^{\overline{K}}_{\un}$ on each coordinate.
\end{conjecture}

\begin{remark}
    Since Chevalley involutions are inner automorphisms for types $A_1,B_n,C_n,D_{2k},E_7$, $E_8,F_4$ and $G_2$, the above conjecture remains the same statement as in \cite{AubertCiubotaruRomano2024}*{Conjecture 9.7}. Thus, Conjecture \ref{conj:main2} holds for $G=\SO_{2n+1}$ by the work of Moeglin--Waldspurger \cite{MoglinWaldspurger2003}, for $G=G_2$ by the work of Ciubotaru--Opdam \cite{CiubotaruOpdam2010} and for $G=\Sp_4$ by the work of Aubert--Ciubotaru--Romano \cite{AubertCiubotaruRomano2024}. 
\end{remark}

In general, the main difficulty of the above conjecture relies on the computation of the unipotent parahoric restriction maps $\res^K_{\un}$. The space $R_{\un}(G)$ has a natural basis in terms of the geometry of the dual group $G^\vee$, while irreducible characters $R_{\un}(\overline{K})$ are parametrized by Lusztig's labels, often of combinatorial flavour. The connection between these parametrizations is deep, and we devote the next section to discuss this in detail for Iwahori-spherical representations. We also give a brief discussion on the computation of this map for non-Iwahori-spherical representations. We closely follow Reeder's discussion in \cite{Reeder2000}*{\S8,9}.

\subsection{Action of the fundamental group}

Let $\mathcal{G}_{\ad}$ be the adjoint group of the simple algebraic $F$-group $\mathcal{G}$ split over $F$ and let $G_{\ad}=\mathcal{G}_{ad}(F)$ be the adjoint group of $F$. The natural short exact sequence of algebraic groups 
$$1\longrightarrow Z(\mathcal{G})\longrightarrow\mathcal{G}\xrightarrow{\ \eta\ }\mathcal{G}_{\ad}\longrightarrow1$$
gives rise, via Galois cohomology, to an exact sequence
$$1\longrightarrow Z(\mathcal{G})(F)\longrightarrow G \xrightarrow{\ \eta\ }G_\ad\xrightarrow{\ \delta\ }H^1(F,Z(\mathcal{G}))\longrightarrow H^1(F,\mathcal{G})=1,$$
by the Kneser--Bruhat--Tits theorem. In particular, $\eta:G\to G_\ad$ is not surjective, and $G_\ad/G:=G_\ad/\eta(G)$ is isomorphic to the Galois cohomology group $H^1(F,Z(\mathcal{G}))$. Let $\mathcal{T}$ be an $F$-split maximal torus of $\mathcal{G}$ and let $\mathcal{T}_\ad:=\mathcal{T}/Z(\mathcal{G})$ be an $F$-split maximal torus of $\mathcal{G}_\ad$, with $T=\mathcal{T}(F)$ and $T_\ad=\mathcal{T}_\ad(F)$. It is known \cite{Springer1998}*{Lemma 16.3.6} that the group $T_\ad/T:=T_\ad/\eta(T)$ is naturally isomorphic to $G_\ad/G$, and this happens to be a very useful characterization.

The standard adjoint action of $G_{\ad}$ on $G$ induces a natural action on $\Irr_{\un}(G)$ given by pullback 
$$(\mathrm{Ad}(g)^*\pi)(h)=(\pi\circ\mathrm{Ad}(g))(h)=\pi(ghg^{-1})\text{ for }g\in G_{\ad}\text{ and } h\in G.$$ 
For any $g\in\eta(G)$, $\mathrm{Ad}(g)$ is an inner automorphism, so the action above factors through $G_\ad/G$ but is generally non-trivial since $G_{\ad}/G$ acts on $G$ by outer automorphisms. In this subsection, we describe the effect of this action on the set of unipotent Langlands parameters of $G$ explicitly. We closely follow Solleveld's results in \cite{Solleveld2020}*{\S2.1} and use Kato's notation in \cite{Kato1983}.

Let $x=su$ be the Jordan decomposition of $x\in G^\vee$ and let $M^\vee:=Z_{G^\vee}(s)^0$ be a connected complex reductive group. We then consider the finite component group
$$R_s:=Z_{G^\vee}(s)/M^\vee=\pi_0(Z_{G^\vee}(s)).$$
Using Galois cohomology, Solleveld constructs a group homomorphism
$$\tau_{s}:G_\ad/G\longrightarrow\Hom(R_s,\CC^\times).$$
Next, we note that $Z_{G^\vee}(x)=Z_{Z_{G^\vee}(s)}(u)$ and therefore $A_x=\pi_0(Z_{Z_{G^\vee}(s)}(u))$. There is a natural map 
$$Z_{G^\vee}(x)\hookrightarrow Z_{G^\vee}(s) \twoheadrightarrow R_s,\quad\text{descending to the map}\quad A_x\xrightarrow{\ \lambda_x\ } R_s,$$
whose kernel is $A^s_u:=\pi_0(Z_{M^\vee}(u))$ and whose image is the subgroup
$$R_s(u)=\left\{r\in R_s\ |\ r\cdot\mathcal{O}_u^{M^\vee}=\mathcal{O}_u^{M^\vee}\right\}.$$
We therefore obtain a (split) short exact sequence
$$1\longrightarrow A_u^s\longrightarrow A_x \xrightarrow{\ \lambda_x\ } R_s(u)\longrightarrow 1,$$
and by composing with $\lambda_x$ we obtain a map
\begin{align}
    \tau_x:G_\ad/G&\longrightarrow \Hom(A_x,\CC^\times),\\
    g &\longmapsto \tau_s(g)\circ\lambda_x.\notag 
\end{align}
We note that, by construction, $\tau_x(g)|_{A^s_u}=1$ for all $g\in G_\ad/G$. By \cite{Solleveld2020}*{Theorem 2.10}, the action of $G_\ad/G$ on the Langlands parameters is given by
$$\mathrm{Ad}(g)^*\pi(x,\phi)=\pi(x,\phi\otimes\tau_x(g)^{-1}).$$
Next, we study in more detail the structure of the group $G_\ad/G$. By the discussion above, it is naturally isomorphic to $T_\ad/T$. Let $\nu_F:F^\times\to\ZZ$ be the normalizer valuation on $F$, and consider the canonical homomorphism
\begin{align*}
    \nu_{T_\ad}:T_\ad&\longrightarrow \Hom(X^*(T_\ad),\ZZ)=X_*(T_\ad)=P^\vee,\\
    t&\longmapsto (\chi\longmapsto\nu_F(\chi(t))),
\end{align*}
where $P^\vee$ is the coweight lattice. The fundamental group of $G_\ad$ is defined as $\Omega:=P^\vee/Q^\vee$, where $Q^\vee=X_*(T)$ is the coroot lattice, and the map $\nu_{T_\ad}$ naturally induces the Kottowitz homomorphism
$$\kappa_T:T_\ad/T \xtwoheadrightarrow{} \Omega$$
fitting in the (split) short exact sequence 
$$1\longrightarrow T_\ad(\mathcal{O}_F)/\eta(T(\mathcal{O}_F))\longrightarrow T_\ad/T\xrightarrow{\ \kappa_T\ }\Omega\longrightarrow 1.$$
The splitting on this short exact sequence is not canonical; it depends on a choice of uniformizer. Moreover, the action of any $t\in T_\ad(\mathcal{O}_F)$ on the Langlands parameters is trivial by \cite{Solleveld2020}*{Lemma 7.5}, and therefore the map $\tau_x$ factors through the Kottowitz homomorphism $\kappa:G_\ad/G\twoheadrightarrow\Omega$, finally giving a map
$$\overline{\tau}_x:\Omega\longrightarrow\Hom(A_x,\CC^\times)$$
such that $\mathrm{Ad}(\omega)^*\pi(x,\phi)=\pi(x,\phi\otimes\overline{\tau}_x(\omega)^{-1})$ for all $\omega\in\Omega$. 

\begin{lemma}\label{lem:action_on_Pi}
    Let $u\in G^\vee$ unipotent and let $(s,h)\in\mathcal{Y}(\Gamma_u)_{\mathrm{ell}}$ be an elliptic pair. For every $\omega\in\Omega$, 
    $$\mathrm{Ad}(\omega)^*\Pi(u,s,h)=\overline{\tau}_{us}(\omega)(h)\Pi(u,s,h).$$
\end{lemma}
\begin{proof}
    This follows by direct computation; indeed
    \begin{align*}
        \mathrm{Ad}(\omega)^*\Pi(u,s,h)&=\sum_{\phi\in\widehat{A_{us}}}\phi(h)\mathrm{Ad}(w)^*\pi(us,\phi)=\\
        &=\sum_{\phi\in\widehat{A_{us}}}\phi(h)\pi(us,\phi\otimes\overline{\tau}_{us}(\omega)^{-1})=\\
        &=\sum_{\psi\in\widehat{A_{us}}}(\psi\otimes\overline{\tau}_{us}(\omega))(h)\pi(us,\psi)=\overline{\tau}_{us}(\omega)(h)\Pi(u,s,h).
    \end{align*} 
\end{proof}

In \cite{Solleveld2020}*{\S2.1}, Solleveld gives a general construction for the homomorphism $\overline{\tau}_x$ for each $x\in G^\vee$. In this setting of unramified Langlands parameters, we can give an explicit formulas that are very useful. Firstly, let $\eta^\vee:G^\vee_{\scn}\to G^\vee$ be the universal covering map of $G^\vee$. Fix some $u\in G^\vee$ unipotent and $(s,h)\in\mathcal{Y}(\Gamma_u)_{\mathrm{ell}}$. While $u$ has a unique unipotent lift $\tilde{u}\in G^\vee_{\scn}$, the semisimple elements $s$ and $h$ do not. Nevertheless, for any lifts $\tilde{s},\tilde{h}\in G^\vee_{\scn}$, the value of
$$c_s(h):=[\tilde{s},\tilde{h}]=\tilde{s}\tilde{h}\tilde{s}^{-1}\tilde{h}^{-1}\in Z(G^\vee_{\scn})$$
is independent of the choice of lifts. Finally, there is a natural identification
$$\Omega\longleftrightarrow X^*(Z(G^\vee_{\scn}))=\Hom(Z(G^\vee_{\scn}),\CC^\times),$$
so each $\omega\in\Omega$ induces a character $\chi_\omega:Z(G^\vee_{\scn})\to\CC^\times$. Then, Solleveld constructs
\begin{equation}
    \overline{\tau}_{us}(\omega)(h)=\chi_\omega(c_s(h)).
\end{equation}
With this expression, we can prove an important identity.
\begin{proposition}\label{prop:identity_tau}
    Fix a unipotent element $u\in G^\vee$ and let $(s,h)\in\mathcal{Y}(\Gamma_u)_{\mathrm{ell}}$ be an elliptic pair. If $C:G^\vee\to G^\vee$ is a Chevalley involution, then for all $\omega\in\Omega$, the identity
    \begin{equation}
        \overline{\tau}_{us}(\omega)(h)=\overline{\tau}_{C(uh)}(\omega)(C(s))
    \end{equation}
    holds.
\end{proposition}
\begin{proof}
    
\end{proof}


Furthermore, the fundamental group $\Omega$ can be naturally embedded inside the automorphism group of the extended Dynkin diagram of $G$. By composing with inner automorphisms if necessary, we may assume that $\mathrm{Ad}(\omega)$ acts on $G$ by permuting simple root spaces according to the automorphism of the Dynkin diagram it is associated with. In particular, if $J\subsetneq S_{\aff}$ and $\omega\in\Omega$, the map $\mathrm{Ad}(\omega)^*$ induces an isomorphism $\overline{K}_J\xrightarrow{\sim} \overline{K}_{\omega(J)}$ and 
\begin{equation}\label{eqn:Omega_action}
    \res_{\un}^{K_J}\mathrm{Ad}(\omega)^*\Pi(u,s,h)\cong\res_{\un}^{K_{\omega(J)}}\Pi(u,s,h).
\end{equation}
\begin{proposition}\label{prop:omega_action}
    Let $G$ be a simple, split $p$-adic group of simply connected type. Suppose that the diagram \eqref{eqn:diagram} commutes for the parahoric subgroup $K_J$ where $J\subsetneq S_\aff$. Then, for any $\omega\in\Omega$, the diagram also commutes for the parahoric $K_{\omega(J)}$.
\end{proposition}
\begin{proof}
    Again, this follows directly from Lemma \ref{lem:action_on_Pi}, Proposition \ref{prop:identity_tau} and \eqref{eqn:Omega_action} since
    \begin{align*}
        \res^{K_{\omega(J)}}_{\un}\circ\FT^\vee_{\mathrm{ell}}\Pi(u,s,h)&=\res^{K_{\omega(J)}}_{\un}\Pi(C(u),C(h),C(s))=\\
        &=\res^{K_J}_{\un}\mathrm{Ad}(\omega)^*\Pi(C(u),C(h),C(s))=\\
        &=\overline{\tau}_{C(uh)}(\omega)(C(s))\ \res^{K_J}_{\un}\circ\FT^\vee_{\mathrm{ell}}\Pi(u,s,h)=\\
        &=\overline{\tau}_{us}(\omega)(h)\ \FT^{K_J}_{\un}\circ\res^{K_J}_{\un}\Pi(u,s,h)=\\
        &=\FT^{K_J}_{\un}\circ\res^{K_J}_{\un}\mathrm{Ad}(\omega)^*\Pi(u,s,h)=\FT^{K_{\omega(J)}}_{\un}\circ\res^{K_{\omega(J)}}_{\un}\Pi(u,s,h).
    \end{align*}
\end{proof}
\iffalse
\begin{remark}\label{rem:condition}
    We note that condition \eqref{eqn:condition} holds when $s=1$ (or $h=1$ by symmetry). Indeed, $C(1)=1$ so $\overline{\tau}_{C(uh)}(\omega)(1)=1$. On the other hand, when $s=1$, $R_s$ is trivial so $A_x=A_u^s$ and therefore $\overline{\tau}_u(\omega)$ are all trivial characters. 
\end{remark}
\fi

\newpage

\iffalse
Let us discuss the image of $\overline{\tau}_x$ in greater detail. Consider the dual map $\eta^\vee:G^\vee_{\scn}\rightarrow G^\vee$ and let $\tilde{u}$ be the unique unipotent lift of $u$ in the simply connected cover $G_{\scn}^\vee$ of $G^\vee$ and fix any lift $\tilde{s}$ of $s$ in $G_{\scn}^\vee$. Then $\tilde{x}=\tilde{s}\tilde{u}$ is a lift of $x$ in $G_{\scn}^\vee$. Since $\tilde{u}$ is the unique unipotent lift of $u$, the canonical map $Z_{G_{\scn}^\vee}(\tilde{u})\rightarrow Z_{G^\vee}(u)$ is surjective. While the map $Z_{G_{\scn}^\vee}(\tilde{s})\rightarrow M^\vee:=Z_{G^\vee}(s)^0$ is well-defined and surjective, the group $Z_{G^\vee}(s)$ may be disconnected, with component group isomorphic to $R$. Consider now the composition
$$Z_{G_{\scn}^\vee}(\tilde{x})=Z_{Z_{G_{\scn}^\vee}(\tilde{s})}(\tilde{u})\twoheadrightarrow Z_{M^\vee}(u)\hookrightarrow Z_{Z_{G^\vee}(s)}(u)=Z_{G^\vee}(x).$$
The first map is surjective (follows from the above observation) while the second map is clearly injective.
Since $Z_{G^\vee}(s)/M^\vee\cong R_s$, we have that $Z_{G^\vee}(x)/Z_{M^\vee}(u)\cong R_s(u)$ by definition. Recall that $A_{\tilde{x}}$ and $A_x$ are the component groups of $Z_{G_{\scn}^\vee}(\tilde{x})$ and $Z_{G^\vee}(x)$, respectively, and let $A_x^0$ be the component group of $Z_{M^\vee}(u)$. It then follows that $A_x\cong A_x^0\rtimes R(u)$, and we have the short exact sequence 
\begin{equation}\label{eqn:SES_comp}
    1\rightarrow Z/(Z\cap Z_{G_{\scn}^\vee}(\tilde{x})^0)\rightarrow A_{\tilde{x}}\rightarrow A_x^0\rightarrow 1,
\end{equation}
which implies that there is a bijection $\phi\leftrightarrow\tilde{\phi}$ between irreducible representations of $A_x^0$ and irreducible representations of $A_{\tilde{x}}$ trivial on the centre $Z$. Since all $\tau_x(g)$ are trivial on $Z_{G^\vee}(s)^0$, it follows that all $\overline{\tau}_x(g)$ are trivial on $A_x^0$ and therefore they all arise as inflations of characters of $R(u)$. In other words, the homomorphism $\overline{\tau}_x(g)$ factors through  
$$\overline{\tau}_x(g):G_{\ad}/G\longrightarrow\Hom(R(u),\CC^\times)\xrightarrow{\mathrm{inf}}\Hom(A_x,\CC^\times).$$
The following result fully determines the image of $\overline{\tau}_x(g)$.
\begin{lemma}
    The map $\overline{\tau}_x:G_{\ad}/G\rightarrow\Hom(R(u),\CC^\times)$ is surjective. 
\end{lemma}
\begin{proof}
    Given the Langlands parameter $(x,\phi)$ of $G$, we let $\tilde{x}\in G^\vee_{\scn}$ be as above and $\tilde{\phi}$ be the character of $A_{\tilde{x}}$ trivial on $Z(G^\vee_{\scn})$ obtained by restricting $\phi$ to $A_x^0$ and the pulling back to $A_{\tilde{x}}$ via \eqref{eqn:SES_comp}. Then $(\tilde{x},\tilde{\phi})$ is naturally a Langlands parameter for $G_{\ad}$. By \cite{Solleveld2020}*{Lemma 7.3}, the pullback 
    $$\eta^*\pi(\tilde{x},\tilde{\phi})=\bigoplus_{\chi\in\widehat{A_x}}\Hom_{A_{\tilde{x}}}(\tilde{\phi},\tilde{\chi})\otimes\pi(x,\chi)=\bigoplus_{\rho\in\widehat{R(u)}}\pi(x,\phi\otimes\rho).$$
    By standard Clifford theory, we know that $G_{\ad}$ must act transitively on the irreducible components of the right hand side, thus completing the proof.
\end{proof}

\newpage
\fi

\subsection{The \texorpdfstring{$\SL_n(F)$}{PDFstring} case} \label{subsec:SLn}

In this section we prove that Conjecture \ref{conj:main2} holds for $G=\SL_n(F)$. Since the dual group is $G^\vee=\PGL_n(\CC)$, via Jordan's normal form each unipotent conjugacy class corresponds to a partition of $n$ $\lambda= (\underbrace{1,\ldots,1}_{r_1}, \underbrace{2,\ldots,2}_{r_2}, \ldots, \underbrace{\ell,\ldots,\ell}_{r_\ell}),$ with $r_1+2r_2+\cdots+\ell r_\ell=n$. Given $u\in\PGL_n(\CC)$ unipotent with corresponding partition $\lambda_u$ of $n$, the reductive part $\Gamma_u$ of the centralizer of $u$ in $G^\vee$ contains elliptic pairs if and only if $\lambda_u$ is rectangular; that is, if $\lambda_u=(\underbrace{m,\ldots,m}_{d})$ for some $n=md$ (see \cite{AubertCiubotaruRomano2024}*{Proposition 13.2}). In that case 
\[
    \PGL_d(\CC)\cong\Gamma_u,\quad g\mapsto g\otimes I_m.
\] 
Let $W=S_d$ be the Weyl group of $\PGL_d(\CC)$ (type $A_{d-1}$) and let $\dot{w}_d$ be the permutation matrix corresponding to the $d$-cycle $w_d=(1,2,3,\ldots,d)\in S_d$. In addition, let $s_d=\mathrm{diag}(1,\zeta_d,\zeta_d^2,\ldots,\zeta_d^{d-1})\in\PGL_d(\CC)$ (note that $s_d$ is a  regular semisimple element). Then a standard group theory argument shows that 
\[
    \mathcal{Y}(\Gamma_u)_{\mathrm{ell}}=\bigsqcup_{k\in(\ZZ/d\ZZ)^\times} \Gamma_u\cdot(s_d,\dot{w}_d^k)
\]
and moreover the pairs $(\dot{w}_d^k,s_d)$ and $(s_d,\dot{w}_d^{-k})$ are $\Gamma_u$-conjugate. In particular, the number of orbits of elliptic pairs in $G^\vee=\PGL_n(\CC)$ is 
\begin{equation*}
    \sum_{u\text{ unipotent class}}|\Gamma_u\backslash\mathcal{Y}(\Gamma_u)_{\mathrm{ell}}|=\sum_{d\mid n}|(\ZZ/d\ZZ)^\times|=\sum_{d\mid n}\varphi(d)=n.
\end{equation*}
For a detailed proof of these facts, see \cite{AubertCiubotaruRomano2024}*{\S 13.1}. 

For the remainder of the subsection, fix some $d\mid n$, let $m=n/d$ and let $u_{d\times m}\in\PGL_n(\CC)$ be the standard unipotent element with Jordan blocks $\lambda_u=(\underbrace{m,\ldots,m}_{d})$. Fix a primitive $d$-th root of unity $\zeta_d$ and let 
$$s_{d\times m}=s_d\otimes I_m=\textrm{diag}(\underbrace{1,\ldots,1}_{m},\underbrace{\zeta_d,\ldots,\zeta_d}_{m},\ldots,\underbrace{\zeta_d^{d-1},\ldots,\zeta_d^{d-1}}_{m})\quad\text{and}\quad \dot{w}_{d\times m}=\dot{w}_d\otimes I_m.$$
%We note that $s_{d\times m}$ and $\dot{w}_{d\times m}$ are the images $s_d$ and $\dot{w}_d$ under the isomorphism $\PGL_d(\CC)\cong\Gamma_u$, $M\mapsto M\otimes I_m$. 
By the discussion above, $(s_{d\times m},\dot{w}_{d\times m}^k)$ is an elliptic pair for $k\in(\ZZ/d\ZZ)^\times$, so we consider the virtual representations 
$$\Pi(u_{d\times m},s_{d\times m},\dot{w}^k_{d\times m})=\sum_{\ell\in\ZZ/d\ZZ}\zeta_d^{k\ell}\pi(x_{d\times m},\phi_\ell),$$
where $x_{d\times m}=u_{d\times m}s_{d\times m}$ and $\phi_l(\dot{w}_{d\times m})=\zeta^\ell_d$. Note that we have used that $A_{x_{d\times m}}=\langle\dot{w}_{d\times m}\rangle\cong\ZZ/d\ZZ$.

The standard Chevalley involution on $\PGL_n(\CC)$ is given by the inverse-transpose map $C(g)=(g^T)^{-1}$, and in particular $C(s_{d\times m})=s_{d\times m}^{-1}$ and $C(w_{d\times m})=w_{d\times m}$. Let $J_m$ be the standard $m$ by $m$ unipotent block in $PGL_m(\CC)$ and let $Y\in\PGL_m(\CC)$ be a matrix such that $Y(J_m^T)^{-1}Y^{-1}=J_m$. Then the matrix 
\begin{equation*}
    X=\begin{pmatrix}
        Y & & & \\
        & & & Y \\
        & & \iddots & \\
        & Y & & 
    \end{pmatrix}\in\PGL_n(\CC)
\end{equation*}
satisfies that $XC(u_{d\times m})X^{-1}=u_{d\times m}$, $Xs_{d\times m}^{-1}X^{-1}=s_{d\times m}$ and $X\dot{w}_{d\times m}X^{-1}=\dot{w}_{d\times m}^{-1}$. Consequently,
\begin{align}\label{eqn:FourierSLn}
    \notag \FT^\vee_{\mathrm{ell}}(\Pi(u_{d\times m}, s_{d\times m},\dot{w}_{d\times m}^{k}))&=\Pi(C(u_{d\times m}), C(\dot{w}_{d\times m}^k), C(s_{d\times m}))\\
    &=\Pi(u_{d\times m}, \dot{w}_{d\times m}^{-k},s_{d\times m})=\Pi(u_{d\times m}, s_{d\times m},\dot{w}_{d\times m}^{k}).
\end{align}

\begin{proposition}
    Conjecture \ref{conj:main2} holds for $G=\SL_n(F)$, without the appearance of roots of unity.
\end{proposition}
\begin{proof}
    According to \eqref{eqn:FourierSLn}, the dual nonabelian Fourier transform $\FT^\vee_{\mathrm{ell}}$ is the identity automorphism on $R_{\un,\mathrm{ell}}(G)$. On the other hand, for each $K\in\max(G)$, the reductive quotient $\overline{K}$ is isomorphic to $\SL_n(\FF_q)$, so $\FT^{\cpt}_{\un}$ is also the trivial automorphism on each $R_{\un}(\overline{K})$. The result follows.
\end{proof}